% Mechanical relocation generated from the preservation ledger.
% Existing prose is included byte-for-byte from preserved blocks.

\section{Comparison of Demand Models}
\label{sec:demand-model-comparison}

The models above answer different questions and are not interchangeable.  The
following comparison makes their principal assumptions and failure modes
explicit.

\begin{center}
\small
\begin{tabular}{>{\raggedright\arraybackslash}p{2.4cm}
                >{\raggedright\arraybackslash}p{3.1cm}
                >{\raggedright\arraybackslash}p{4.8cm}
                >{\raggedright\arraybackslash}p{5.0cm}}
\hline
Model & Unknown scale & Principal identifying information or assumption & Main limitation and role \\
\hline
Cell deviations
& One parameter per free OD-time cell
& Detailed count response, baseline matrix, and cell-level prior or penalty
& Flexible but intrinsically weakly identified and high dimensional; use for small problems or targeted refinement. \\

Conditional gravity
& A few shared coefficients, optionally plus a small production basis
& Feasible alternatives, journey productions, attractiveness, and common responses to travel time and transfers
& Misspecification can propagate systematically across many cells; preferred full-network model with MAP and diagnostics. \\

Route-level IPF
& One triangular vehicle-leg table per route pattern
& Observed route boarding and alighting marginals plus a positive seed
& Does not identify complete journeys or distinguish transfers; audit baseline and initializer only. \\

Entropic transport
& One flow per supported OD cell, implicitly selected by entropy
& Genuine journey marginals, generalized costs, support, and regularization or marginal penalties
& Solution depends on supplied marginals and entropy choices; auxiliary baseline or initializer. \\
\hline
\end{tabular}
\end{center}

\subsection{Consequences for Interpretation}

Parameter count is not the same as information content.  Cell-level estimation
has the largest parameter space but does not become more informative merely by
being flexible.  Gravity reduces dimension by asserting that many OD cells
share behavioral coefficients and conditional normalization.  IPF and entropy
obtain determinate numerical tables by imposing marginal and seed or entropy
structure; this determinacy must not be reported as identification from stop
counts.

Likewise, the models estimate different objects.  Route-level IPF estimates
vehicle-leg movements.  The other models target passenger journeys, but
entropy requires externally meaningful journey marginals, while conditional
gravity requires journey productions or a declared model for them.  Raw APC
boardings include transfer boardings and therefore cannot silently substitute
for either journey origins or gravity productions.

\subsection{Recommended Modeling Hierarchy}

For a large network, the primary specification is the smallest conditional
gravity model that is scientifically defensible, estimated by MAP.  Its low
dimension addresses both computation and the intrinsic underdetermination of
cell-level OD recovery.  Increasing model complexity should be evidence-led:
first examine residual structure, curvature, grouped holdout, and sensitivity;
then introduce a narrowly defined temporal, spatial, production, or transfer
effect and verify that it improves adequacy without destroying identification.

Cell-level deviations remain a controlled comparison and local refinement
tool.  Route-level IPF and entropic transport supply audits or initial values,
not alternative evidence that overrides the journey model.  Estimates from
different rows of the table should therefore be compared at common observable
quantities---especially predicted boarding and alighting counts---and not only
by comparing OD matrices whose meanings and identifying assumptions differ.
